\section{Problem 252: an explicit phase obstruction for the fourth divisor power}
\subsection*{1) OBJECTIVE AND SOURCE REVIEW}
Attempt dated 2026-09-23. The target is irrationality of $S_k=\sum_{n\ge1}\sigma_k(n)/n!$ for every fixed positive integer $k$. I read \texttt{252.tex} and \texttt{252\_v2.tex} fully, including their tail estimates and conditional prime patterns. A literature correction matters: Schlage-Puchta's cited paper already proves $k=3$ unconditionally, as well as the all-$k$ implication from Schinzel's hypothesis H. The following is a worked conditional obstruction for $k=4$, not a new proof of that known conditional theorem.
\subsection*{2) WORK}
Put $D=24^4$. Suppose there are arbitrarily large $q\equiv1\pmod D$ for which all four numbers
\[
q,\quad(q+1)/2,\quad(q+2)/3,\quad(q+3)/4
\]
are prime. Assume toward contradiction that $S_4$ is rational. For large $q$, multiplication by $(q-1)!$ shows that
\[
T_q=\sum_{i\ge1}\frac{\sigma_4(q+i-1)}{q(q+1)\cdots(q+i-1)}
\]
is an integer. The terms $i\ge5$ total $O(1/q)$: use $\sigma_4(m)\le\zeta(4)m^4$ and factorial decay. The first term is $q^3+1/q$. For $i=2,3,4$, the prime $r=(q+i-1)/i$ exceeds $i$, so multiplicativity gives $\sigma_4(ir)=\sigma_4(i)(r^4+1)$. Polynomial division consequently yields
\[
T_q=q^3+\frac{17}{16}(q^2+3q+3)
+\frac{82}{81}(q+5)+\frac{273}{256}+O(1/q).
\]
All three rational-polynomial values modulo one depend only on $q\bmod D$. At $q=1$ their sum is
\[
\frac{119}{16}+\frac{164}{27}+\frac{273}{256}
=\frac{100763}{6912}\equiv\frac{3995}{6912}\pmod1.
\]
This nonzero fixed phase contradicts integrality for sufficiently large $q$. Thus the specified prime-pattern infinitude would imply irrationality of $S_4$.

There is no elementary local obstruction to the pattern: writing $q=Dt+1$, the four forms are $(D/i)t+1$. Modulo $2$ or $3$ each is $1$; modulo every prime $\ell\ge5$ they exclude at most four residues, fewer than $\ell$.
\subsection*{3) ADVERSARIAL VERIFICATION}
Local admissibility is not a proof of simultaneous primality. The argument's missing infinitude is exactly an unproved prime-pattern assertion. The polynomial remainders are $O(1/q)$, including the additional $+1$ in each prime divisor sum; they cannot cancel the fixed phase. Nothing here proves the $k=4$ case unconditionally, and the already established $k=3$ case must not be advertised as new progress.
\subsection*{4) FINAL}
\textbf{UNRESOLVED.} An explicit rationality obstruction replaces an unspecified phase, but it still requires an unproved simultaneous-prime input. The full all-$k$ question is not settled.
\subsection*{5) SOURCES CHECKED}
Both complete supplied versions; J.-C. Schlage-Puchta, \emph{The irrationality of a number theoretical series}, \texttt{https://arxiv.org/abs/1105.1452}, checked 2026-09-23. The displayed divisions and phase are exact elementary calculations.
\subsection*{6) COMPLETION ESTIMATE}
Subjective progress toward an unconditional solution for every $k$.
\textbf{COMPLETION: 15\%}
